%% =================================================================
%% Beber, Lamon, Saveriano, Fontanelli, Palopoli.
%% "Autonomous Robotic Tissue Palpation and Abnormalities
%%  Characterisation via Ergodic Exploration."
%% IEEE Robotics and Automation Letters, preprint accepted Feb. 2026.
%% DOI: 10.1109/LRA.2026.3673907
%%
%% Reconstruction of page 1 as study notes.
%% Prose: extracted verbatim from the PDF text layer via pdftotext
%%        (column interleaving unbraided, right-column figure labels
%%        discarded as non-prose).
%% Layout: hand-styled single-column (not IEEE two-column — study
%%        notes, not pixel-perfect clone). Figure 1 as placeholder box.
%% =================================================================

\documentclass[11pt]{article}

\usepackage[margin=1in]{geometry}
\usepackage{amsmath, amssymb}
\usepackage{mathrsfs}
\usepackage{xcolor}
\usepackage{fancyhdr}
\usepackage{url}

\pagestyle{fancy}
\fancyhf{}
\lhead{\small L.~Beber et al.}
\rhead{\small IEEE RA-L Preprint, Accepted Feb.\ 2026}
\cfoot{\small 1 $\to$ \arabic{page}}
\renewcommand{\headrulewidth}{0.4pt}

\setcounter{page}{1}
\setlength{\parskip}{4pt}

\begin{document}

% ---------- Preprint notice strip (top of page) ----------
\begin{center}
{\footnotesize\itshape
This article has been accepted for publication in IEEE Robotics and
Automation Letters. This is the author's version which has not been
fully edited and content may change prior to final publication.
Citation information: DOI 10.1109/LRA.2026.3673907\par}
\end{center}

\vspace{0.4em}
\hrule
\vspace{0.6em}

\noindent
{\footnotesize\scshape IEEE Robotics and Automation Letters.\ Preprint Version.\ Accepted February, 2026 \hfill 1}

\vspace{1.2em}

% ---------- Title ----------
\begin{center}
{\LARGE\bfseries
Autonomous Robotic Tissue Palpation and Abnormalities\\
Characterisation via Ergodic Exploration\par}
\end{center}

\vspace{0.6em}

% ---------- Authors ----------
\begin{center}
Luca Beber\textsuperscript{1},
Edoardo Lamon\textsuperscript{1},
Matteo Saveriano\textsuperscript{2},
Daniele Fontanelli\textsuperscript{2},
and Luigi Palopoli\textsuperscript{1}
\end{center}

\vspace{1em}

% ---------- Abstract ----------
\noindent
\textbf{\textit{Abstract}---}We propose a novel autonomous robotic
palpation framework for real-time elastic mapping during tissue
exploration using a viscoelastic tissue model. The method combines
force-based parameter estimation using a commercial force/torque
sensor with an ergodic control strategy driven by a tailored Expected
Information Density, which explicitly biases exploration toward
diagnostically relevant regions by jointly considering model
uncertainty, stiffness magnitude, and spatial gradients. An Extended
Kalman Filter is employed to estimate viscoelastic model parameters
online, while Gaussian Process Regression provides spatial modelling
of the estimated elasticity, and a Heat Equation Driven Area Coverage
controller enables adaptive, continuous trajectory planning.
Simulations on synthetic stiffness maps demonstrate that the proposed
approach achieves better reconstruction accuracy, enhanced
segmentation capability, and improved robustness in detecting stiff
inclusions compared to Bayesian Optimisation-based techniques.
Experimental validation on a silicone phantom with embedded inclusions
emulating pathological tissue regions further corroborates the
potential of the method for autonomous tissue characterisation in
diagnostic and screening applications.

\vspace{0.4em}

\noindent
\textit{Index Terms}---Medical Robots and Systems; Sensor-based
Control; Force and Tactile Sensing

\vspace{1em}

% ---------- Figure 1 placeholder ----------
\begin{center}
\fbox{\parbox{0.92\textwidth}{\centering\vspace{1em}
\textbf{[Figure 1 --- placeholder]}\\[0.8em]
Photograph of the experimental setup, arranged as a composite of
three panels. \textbf{Bottom-left}: top-down view of a silicone
phantom sample --- a roughly square slab with a visible darker
``Abnormality'' inclusion near its centre. \textbf{Top-left}: the
corresponding ground-truth elasticity map obtained via a dense grid
scan, shown as a colour-coded 2D field with a vertical colour bar
labelled \textit{Elastic Modulus (kPa)}. \textbf{Right}: a Universal
Robots UR3e manipulator shown palpating the silicone sample, with
labelled call-outs identifying the \textit{F/T Sensor} mounted at the
wrist and the cylindrical \textit{Indenter} tip in contact with the
sample surface.
\vspace{1em}}}
\end{center}

\vspace{0.3em}

\begin{center}
\textbf{Fig. 1.} Experimental setup for the proposed ergodic search:
(bottom-left) silicone sample; (top-left) ground-truth elasticity map
obtained using a grid-based approach; (right) the robotic manipulator
palpating the silicone sample.
\end{center}

\vspace{1.2em}

% ---------- Introduction ----------
\begin{center}
\textsc{I.\ Introduction}
\end{center}

\noindent
{\large\bfseries P}ALPATION is a fundamental clinical procedure that
allows physicians, by touch, to assess tissue stiffness, texture, and
the presence of abnormalities such as tumours or swollen lymph nodes
[1]. Its speed and non-invasiveness make it crucial for early
diagnosis, particularly of conditions such as breast cancer. However,
palpation is inherently subjective and sensitive to inter-practitioner
variability, which may compromise diagnostic accuracy and
repeatability for small or deep-seated lesions. For this reason,
physical examinations are often complemented or replaced by imaging
techniques such as mammography, the gold standard for early detection
[2], and ultrasound [3].

Accurate assessment of tissue mechanics is equally critical in
robot-assisted minimally invasive surgery (RMIS), where robots must
adapt to variable stiffness for safe navigation and manipulation,
especially for autonomous systems operating with minimal human
supervision [4]. To support such needs, several technologies have
emerged. Imaging-based elastography methods, e.g., ultrasound [5],
magnetic resonance [6], or laser-based methods [7], estimate tissue
stiffness but typically provide relative measures, are costly or slow,
and cannot be easily integrated into RMIS. Tactile sensing, in
contrast, offers real-time, quantitative stiffness measurements by
capturing force and deformation. Previous works have shown its
potential to reduce applied forces, shorten procedures, and improve
surgical outcomes [8], [9]. Early stiffness mapping used uniform
probing grids [10], [11], while later methods adopted Bayesian
Optimisation (BO) to reduce the number of probing actions [12]--[14],
often coupled with custom tactile sensors [8], [9], [15], guided by
anatomical priors [16] or with ad-hoc refined segmentation strategies
[15], [17].

In this work, we propose a method for autonomous robotic palpation
that enables continuous estimation of elastic tissue stiffness using a
viscoelastic contact model, relying solely on a commercial
force/torque (F/T) sensor. The use of a viscoelastic model allows more
accurate representation of tissue dynamics during palpation compared
to purely elastic formulations, while abnormality characterisation in
this work focuses on the estimated elastic stiffness. This sensing
strategy provides a practical alternative to methods based on embedded
or specialised sensors, which typically require customised hardware
integration and lack compatibility with standard robotic platforms. In
contrast, the use of an off-the-

% ---------- Bottom-of-page footnote block ----------
\vfill
\hrule
\vspace{0.4em}
\noindent
{\footnotesize
Manuscript received: October, 15, 2025; Revised December, 16, 2025;
Accepted February, 23, 2026.\\[2pt]
This paper was recommended for publication by Editor Burgner-Kahrs
Jessica upon evaluation of the Associate Editor and Reviewers'
comments. This research has been funded by the the PNRR project FAIR
- Future AI Research (PE00000013), by the European Union projects
INVERSE (GA no. 101136067) and MAGICIAN (GA no. 101120731) and by the
MUR Departments of Excellence 2023-27 program (L.232/2016).\\[2pt]
\textsuperscript{1}Department of Information Engineering and Computer
Science, University of Trento, Trento, Italy
\texttt{name.surname@unitn.it}\\[2pt]
\textsuperscript{2}, Department of Industrial Engineering, University
of Trento, Trento, Italy \texttt{name.surname@unitn.org}\\[2pt]
Digital Object Identifier (DOI): see top of this page.\par}

\vspace{0.4em}
\hrule
\vspace{0.3em}

\begin{center}
{\footnotesize\itshape
This work is licensed under a Creative Commons Attribution 4.0
License. For more information, see
\url{https://creativecommons.org/licenses/by/4.0/}\par}
\end{center}

\end{document}
